\documentclass[a4paper,11pt]{article}
\usepackage[utf8]{inputenc}
\usepackage[english]{babel} % Adds processing for some simple special characters.

% Packages for better page formatting and nice footnotes. %

\usepackage{fancyhdr} % Headers.
\usepackage{multicol} % Columns.
\usepackage[margin=0.75in]{geometry} % Margins.

\pagestyle{fancy}
\fancyhf{} % What does this do?
\lhead{Alexander S. Wheaton}
\rhead{\thedate}
\cfoot{Page \thepage}

\renewcommand{\thefootnote}{\fnsymbol{footnote}}

% Title and author for title page. %

\title{Deep Stacking of AGN Hypervariables}
\author{
    Alexander S. Wheaton\\
    School of Physics and Astronomy\\
    The University of Edinburgh\\
    s1572046@ed.ac.uk\break
}

% Utilities for fine-grained control over image insertion. %

\usepackage{graphicx} % Insert images.
\usepackage{float} % Images float in document environment.
\usepackage{wrapfig} % Image/tables in multicols with \wrapfigure or \wraptable.
\usepackage{caption} % Captions for single images.
\usepackage{subcaption} % Captions for simultaneous images.

% Not sure what this is for. %

\usepackage{fullpage,epsf}
\usepackage{lipsum}

% Some utilities for scientific and mathematical writing. %

\usepackage{siunitx} % Formatting for numbers with SI units.
\usepackage{mathabx} % Astronomical symbols.
\usepackage{isotope} % Nice markup syntax for chemical symbols.
\usepackage{xfrac}   % Slant fractions and other utilities.

% Utilities for manual kerning adjustment. %

\newcommand\K{\kern.05em} % Small amount of kerning.
\newcommand\KK{\kern.1em} % Medium amount of kerning.
\newcommand\KKK{\kern.2em} % Large amount of kerning.
\newcommand\KKKK{\kern.3em} % Very large amount of kerning.

% Python code-highlighting syntax. %

\usepackage{listings}
\usepackage{color}
\definecolor{codegreen}{rgb}{0,0.6,0}
\definecolor{codegray}{rgb}{0.5,0.5,0.5}
\definecolor{codepurple}{rgb}{0.58,0,0.82}
\definecolor{backcolour}{rgb}{0.95,0.95,0.92}
\definecolor{deepblue}{rgb}{0,0,0.5}
\definecolor{deepred}{rgb}{0.6,0,0}
\definecolor{deepgreen}{rgb}{0,0.5,0}
\lstdefinestyle{python}{
    backgroundcolor=\color{backcolour},
    commentstyle=\color{codegreen},
    keywordstyle=\color{deepblue},
    numberstyle=\tiny\color{codegray},
    stringstyle=\color{codegreen},
    basicstyle=\footnotesize,
    breakatwhitespace=false,
    breaklines=true,
    captionpos=b,
    keepspaces=true,
    numbers=left,
    numbersep=5pt,
    showspaces=false,
    showstringspaces=false,
    showtabs=false,
    tabsize=2
}

\lstset{style=python}

% Bibliography and referencing style.

\usepackage[backend=bibtex,style=phys,sorting=ynt]{biblatex} % Use style=draft for troubleshooting.
\addbibresource{references.bib}

% Page formatting.

\begin{document}

\pagestyle{empty}                       % No numbers of title page
\epsfxsize=40mm                         % Size of crest
\begin{minipage}[b]{110mm}
    {\Huge\bf School of Physics\\ and Astronomy
    \vspace*{17mm}}
\end{minipage}
\hfill
\begin{minipage}[t]{40mm}
    \makebox[40mm]{
    \includegraphics[width=40mm]{img/crest.eps}}
\end{minipage}
\par\noindent                                           % Centre Title, and name
\vspace*{2cm}
\begin{center}
    \Large\bf \Large\bf Senior Honours Project\\
    \Large\bf Astrophysics\\[10pt]                     % Change to MP/CP/Astro
    \LARGE\bf Deep Stacking of AGN Hypervariables
\end{center}
\vspace*{0.5cm}
\begin{center}
    \bf Alexander S. Wheaton\\
    \today
\end{center}
\vspace*{5mm}

\begin{abstract}
        \lipsum[1]
\end{abstract}

\vspace*{1cm}

\subsubsection*{Declaration}
\begin{quotation}
        I declare that this project and report is my own work.
\end{quotation}

%\vspace*{2cm}
Signature:\hspace*{8cm}Date: \today

\vfill
{\bf Supervisor:} Professor A. Lawrence, FRSE, FRaS
\hfill
12 Weeks

\newpage
\setcounter{page}{1} % Set page number to 1
\tableofcontents
\newpage
\section{Introduction}
\subsection{AGNs and the Viscous Accretion Disc}
\cite{lawrence_2018}
\subsection{The Pan-STARRS Survey}
\cite{bruce_2017}

\section{The Liverpool Telescope Data}
\subsection{S/N Ratio and Astronomical Seeing}
\section{Centroiding on J094511}
\subsection{Initial Guess: Max-Value Centroiding}
\subsection{Initial Guess: WCS Centroiding}
\cite{lang_2010}
\cite{greisen_2002}
\cite{calabretta_2002}
\subsection{Refinement: Two Dimensional Weighted Mean}
\subsection{Aperture Centroiding}
\subsection{Challenges in the U-band}
\section{Cross Sections of J094511}
\subsection{Linear and Radial Cross Sections}
\subsection{Excess Luminosity in the Wings}
\subsection{Point Spread Function Subtraction}

\section{Conclusion}

\printbibliography

\section*{Appendix A: Alignment Algorithm}

\end{document}
